% ────────────────────────────────────────────
%  Masterprüfung — Präsentation (TU Graz)
%  Literatur kommt aus derselben Bibliothek wie die Masterarbeit.
% ────────────────────────────────────────────
\documentclass[aspectratio=169, 11pt]{beamer}
\usetheme{TUGraz}

\usepackage[utf8]{inputenc}
\usepackage{csquotes}
\usepackage{booktabs}
\usepackage{amsmath}
\usepackage{pgfplots}\pgfplotsset{compat=1.18}
% Literatur: biblatex + biber (natbib-Befehle wie \citep/\citet funktionieren weiter)
\usepackage[backend=biber, style=authoryear, natbib=true, maxcitenames=2, uniquename=false]{biblatex}
\addbibresource{references.bib}

\title{Hopfen, Hefe, Hochgenuss}
\subtitle{Gärtemperatur und Aroma obergäriger Biere -- Masterprüfung}
\author{Norbert Winter}
\institute{Institut für Brau- und Gärungstechnologie (fiktiv) \textbullet{} TU Graz}
\date{\today}

\begin{document}

\begin{frame}[plain]
  \titlepage
\end{frame}

\begin{frame}{Agenda}
  \tableofcontents
\end{frame}

\section{Motivation}

\begin{frame}{Warum Bier?}
  \begin{itemize}
    \item Nur vier Zutaten -- und doch tausende Aromen
    \item Ein Großteil des Aromas entsteht im \alert{Gärtank}
    \item Systematische Daten zur Gärtemperatur fehlen \citep{malz2021ester}
  \end{itemize}
  \vspace{4mm}
  \begin{block}{Forschungsfrage}
    Wie stark prägt die Gärtemperatur das Aromaprofil obergäriger Biere?
  \end{block}
\end{frame}

\section{Methodik}

\begin{frame}{Vorgehen}
  \begin{columns}[T]
    \begin{column}{0.48\textwidth}
      \begin{enumerate}
        \item 12 Sude Pale Ale
        \item 4 Temperaturstufen (16--24\,°C)
        \item Gaschromatografie
        \item Blindverkostung ($n=24$)
      \end{enumerate}
    \end{column}
    \begin{column}{0.48\textwidth}
      \begin{exampleblock}{Arrhenius-Ansatz}
        \[ k(T) = A \cdot \exp\!\left(-\frac{E_a}{R\,T}\right) \]
      \end{exampleblock}
    \end{column}
  \end{columns}
\end{frame}

\section{Ergebnisse}

\begin{frame}{Esterbildung}
  \centering
  \begin{tikzpicture}
    \begin{axis}[width=0.78\textwidth, height=5.4cm, xlabel={Gärtemperatur [°C]}, ylabel={mg/L},
      grid=major, grid style={black!8}, legend pos=north west, legend style={draw=none, font=\small, fill=none},
      every axis plot/.append style={very thick}]
      \addplot[color=tured, mark=*] coordinates {(16,1.1) (18,1.5) (20,1.9) (22,2.6) (24,3.1)};
      \addplot[color=tugray, mark=square*] coordinates {(16,0.4) (18,0.5) (20,0.6) (22,0.7) (24,0.8)};
      \legend{Isoamylacetat (Banane), {Ethylhexanoat (Apfel, ×2)}}
    \end{axis}
  \end{tikzpicture}
\end{frame}

\begin{frame}{Verkostung}
  \centering
  \begin{tabular}{lrrr}
    \toprule
    Gärtemperatur & Fruchtigkeit & Sauberkeit & Gesamt \\
    \midrule
    16\,°C & 4{,}1 & 7{,}8 & 6{,}2 \\
    \alert{18\,°C} & \alert{5{,}6} & \alert{7{,}2} & \alert{7{,}4} \\
    20\,°C & 6{,}9 & 5{,}9 & 6{,}8 \\
    24\,°C & 8{,}2 & 3{,}7 & 5{,}1 \\
    \bottomrule
  \end{tabular}
\end{frame}

\section{Fazit}

\begin{frame}{Fazit \& Ausblick}
  \begin{itemize}
    \item<1-> +4\,K $\Rightarrow$ rund +38\,\% Isoamylacetat
    \item<2-> 18\,°C ist der Sweet Spot für ein Pale Ale
    \item<3-> Nächster Schritt: Kveik-Hefen und Temperaturprofile
  \end{itemize}
\end{frame}

\danke

% ── Backup / Anhang ──
\begin{frame}{Literatur}
  \scriptsize
  \printbibliography[heading=none]
\end{frame}

\end{document}
